\section{Results and Analysis}
\label{sec:results}

This section reports the main cross-day results, controlled ablations, deployment-shift evaluation, edge-deployment benchmark, and cross-receiver stress test. We emphasize file-level accuracy and file-level macro-F1 because they directly correspond to capture-file-level device authentication decisions.

\subsection{Cross-Day Main Results}

Table~\ref{tab:cross_day_main} summarizes the clean cross-day results under the Day1--3/Day4/Day5 protocol. The reproduced CNN-IQ baseline achieves 54.2$\pm$14.2\% file-level accuracy and 45.6$\pm$14.8\% file-level macro-F1. The proposed cross-attentive OOB hybrid reaches 75.0$\pm$5.3\% file-level accuracy and 67.9$\pm$6.8\% file-level macro-F1, improving CNN-IQ by 20.8 percentage points in file-level accuracy.

\begin{table}[!t]
\centering
\caption{Cross-day main results on Day5 test (mean $\pm$ std over seeds, \%).}
\label{tab:cross_day_main}
\footnotesize
\setlength{\tabcolsep}{2.5pt}
\renewcommand{\arraystretch}{1.05}
\resizebox{0.98\columnwidth}{!}{%
\begin{tabular}{@{}lcccc@{}}
\toprule
Model & F-Acc & F-F1 & W-Acc & W-F1 \\
\midrule
CNN-IQ & $54.2{\pm}14.2$ & $45.6{\pm}14.8$ & $43.5{\pm}5.2$ & $38.7{\pm}5.7$ \\
RF-HSTU no OOB & $66.7{\pm}3.4$ & $59.1{\pm}4.5$ & $39.4{\pm}13.7$ & $36.8{\pm}14.3$ \\
\textbf{F cross-attn + chirp} & $\mathbf{75.0{\pm}5.3}$ & $\mathbf{67.9{\pm}6.8}$ & $41.5{\pm}2.4$ & $39.4{\pm}2.1$ \\
\bottomrule
\end{tabular}%
}
\end{table}


The improvement is not only reflected in the mean value but also in seed-level stability. As shown in Fig.~\ref{fig:cross_day_seed_bars}, the proposed model outperforms CNN-IQ in four out of five random seeds and ties in the remaining seed. A pooled paired bootstrap analysis over file-level predictions yields a 95\% confidence interval of [+9.2, +32.5] percentage points for the accuracy gain over CNN-IQ. This interval is entirely above zero, indicating that the improvement is unlikely to be explained by seed selection alone.

\begin{figure}[!t]
    \centering
    \includegraphics[width=\linewidth]{fig2_cross_day_seed_bars.pdf}
    \caption{Cross-day file-level accuracy across five random seeds. The proposed cross-attentive OOB hybrid improves over CNN-IQ in four seeds and ties in one seed.}
    \label{fig:cross_day_seed_bars}
\end{figure}

The diagnostic RF-HSTU variant without OOB evidence also improves over CNN-IQ in file-level accuracy, reaching 66.7$\pm$3.4\%. This suggests that RF patch sequence modeling itself is useful for cross-day generalization. However, the final OOB-guided cross-attentive model further improves accuracy and macro-F1, showing that OOB spectral evidence provides additional discriminative information when fused appropriately.

\subsection{Fusion and Chirp Ablation}

Table~\ref{tab:fusion_chirp} reports the controlled fusion/chirp ablation. The goal is to separate the effect of the OOB fusion mechanism from the effect of chirp-aware embeddings. The results show a clear pattern: cross-attentive OOB fusion is the key factor behind stable cross-day performance, while chirp embedding alone does not explain the improvement.

\begin{table}[!t]
\caption{Fusion/Chirp Ablation under Cross-Day Protocol}
\label{tab:fusion_chirp}
\centering
\renewcommand{\arraystretch}{1.1}
\begin{tabular}{lccccc}
\toprule
Fusion & Chirp & File-Acc & File-Macro-F1 & $n$ seeds \\
\midrule
concat & $\times$ & 18.3$\pm$18.6 & 12.3$\pm$15.7 & 5 \\
concat & \checkmark & 9.7$\pm$2.0 & 3.1$\pm$1.8 & 3 \\
cross-attn & $\times$ & 75.0$\pm$3.4 & 68.1$\pm$4.0 & 3 \\
cross-attn & \checkmark & \textbf{75.0$\pm$5.3} & \textbf{67.9$\pm$6.8} & 5 \\
\bottomrule
\end{tabular}
\end{table}


The cross-attentive model without chirp embedding achieves 75.0$\pm$3.4\% file-level accuracy, essentially matching the full model with chirp embedding at 75.0$\pm$5.3\%. In contrast, the concat OOB model remains unstable. Without chirp embedding, concat OOB reaches only 18.3$\pm$18.6\%, and adding chirp embedding further results in 9.7$\pm$2.0\%. Therefore, chirp embedding does not rescue naive concatenation.

\begin{figure}[!t]
    \centering
    \includegraphics[width=\linewidth]{fig3_fusion_chirp_ablation.pdf}
    \caption{Fusion/chirp ablation. Cross-attentive OOB fusion remains stable with or without chirp embedding, while concat-based OOB fusion remains weak and unstable.}
    \label{fig:fusion_chirp_ablation}
\end{figure}

Prediction-distribution diagnostics further support this interpretation. The cross-attentive variants produce predictions over many device classes, while concat-based variants often collapse to a small number of dominant classes. This indicates that the failure of concat OOB fusion is not simply a lack of chirp prior, but a more fundamental instability in how OOB evidence is injected into the RF representation. We therefore treat cross-attentive OOB fusion as the primary mechanism of the proposed model, and chirp embedding as a compatible auxiliary LoRa-structure prior.

\subsection{Deployment-Shift Evaluation}

Table~\ref{tab:deployment_shift} reports deployment-shift evaluation under configuration, location, and distance leave-one-domain-out protocols. These results are used as extended deployment evidence rather than as the primary clean cross-day protocol, as their training recipe follows the Phase4 setting described in Section~\ref{sec:experiments}.

\begin{table}[!t]
\centering
\caption{Deployment-shift LOCO results (file-level accuracy, \%).}
\label{tab:deployment_shift}
\scriptsize
\setlength{\tabcolsep}{2.1pt}
\renewcommand{\arraystretch}{1.04}
\begin{adjustbox}{max width=\columnwidth,center}
\begin{tabular}{@{}llcccl@{}}
\toprule
Shift & Held-out & CNN & F & Win & Note \\
\midrule
Config & Config1 (SF7 BW125) & 0.0 & 4.2 & F & both weak \\
Config & Config2 (SF8) & 4.2 & 8.3 & F & \\
Config & Config3 (SF11) & 8.3 & 16.7 & F & \\
Config & Config4 & 4.2 & 0.0 & CNN & \\
Location & room & 29.2 & 54.2 & F & indoor \\
Location & office & 37.5 & 54.2 & F & \\
Location & outdoor & 25.0 & 16.7 & CNN & difficult \\
Distance & 5 m & 8.3 & 8.3 & tie & \\
Distance & 10 m & 4.2 & 20.8 & F & \\
Distance & 15 m & 4.2 & 33.3 & F & \\
Distance & 20 m & 8.3 & 37.5 & F & \\
\bottomrule
\end{tabular}
\end{adjustbox}
\end{table}


The distance-shift results provide the clearest deployment benefit. At 10\,m, 15\,m, and 20\,m held-out distances, the hybrid model achieves 20.8\%, 33.3\%, and 37.5\% file-level accuracy, respectively, compared with 4.2\%, 4.2\%, and 8.3\% for CNN-IQ. Fig.~\ref{fig:distance_shift} visualizes this trend. The advantage becomes more visible at larger distances, suggesting that OOB-aware RF modeling can preserve device evidence when propagation distance changes.

\begin{figure}[!t]
    \centering
    \includegraphics[width=\linewidth]{fig4_distance_shift.pdf}
    \caption{Distance leave-one-domain-out evaluation. The hybrid model shows stronger performance than CNN-IQ at 10\,m, 15\,m, and 20\,m held-out distances.}
    \label{fig:distance_shift}
\end{figure}

The location-shift results are more mixed. The hybrid model improves over CNN-IQ in indoor and office held-out locations, but underperforms CNN-IQ in the outdoor held-out location. Configuration shift is challenging for both models, with all file-level accuracies remaining low. These results indicate that the proposed model improves robustness under some realistic deployment changes, especially distance and indoor-location shifts, but does not solve all forms of deployment variability.

\subsection{Edge Deployment Benchmark}

Table~\ref{tab:edge_deployment} summarizes the edge-deployment benchmark. CNN-IQ is substantially smaller, with 47.7K parameters, while the proposed hybrid model has approximately 1.16M parameters. However, the latency overhead remains moderate in the evaluated setting: batch-1 latency increases from approximately 1.12\,ms for CNN-IQ to 2.45\,ms for the hybrid model.

\begin{table}[!t]
\caption{Edge Deployment Benchmark (Phase6)}
\label{tab:edge_deployment}
\centering
\renewcommand{\arraystretch}{1.1}
\begin{tabular}{lcc}
\toprule
Metric & CNN-IQ & Hybrid \\
\midrule
Parameters & 47.7K & 1.16M \\
Latency @ bs=1 (ms) & 1.12 & 2.45 \\
Latency @ bs=32 (ms) & 1.10 & 2.53 \\
Peak GPU mem @ bs=32 (MB) & 66.4 & 123.6 \\
\bottomrule
\end{tabular}
\end{table}


These results show that the proposed model is more computationally expensive than the CNN baseline, but still feasible for gateway-side inference. Since LoRa RFFI authentication is performed at the edge gateway rather than on battery-constrained transmitters, the additional computation is acceptable for many deployment scenarios. We therefore view the hybrid model as a practical gateway-side authentication module, while further compression and hardware-specific optimization are left for future work.

\subsection{Cross-Receiver Stress Test}

Finally, Table~\ref{tab:cross_receiver} reports strict source-only cross-receiver transfer. This evaluation is intentionally difficult: the model is trained and validated on one receiver and tested on the other receiver, with no target-receiver labels used for checkpoint selection.

\begin{table}[!t]
\caption{Cross-Receiver Stress Test (Phase5-Clean, Strict Source-Only)}
\label{tab:cross_receiver}
\centering
\renewcommand{\arraystretch}{1.1}
\begin{tabular}{lcc}
\toprule
Direction & CNN-IQ & F Hybrid \\
\midrule
RX1$\rightarrow$RX2 & 4.2$\pm$0.0 & 18.1$\pm$3.9 \\
RX2$\rightarrow$RX1 & 23.6$\pm$15.3 & 15.3$\pm$7.1 \\
\bottomrule
\end{tabular}
\end{table}


The results are direction-dependent. In RX1$\rightarrow$RX2 transfer, CNN-IQ remains at 4.2$\pm$0.0\%, which is near the random baseline of 4.17\%, while the proposed hybrid model improves to 18.1$\pm$3.9\%. This suggests that OOB-guided fusion can provide limited benefit in this transfer direction. However, in RX2$\rightarrow$RX1 transfer, CNN-IQ obtains 23.6$\pm$15.3\%, whereas the hybrid model reaches 15.3$\pm$7.1\%. The large seed variance and direction dependence show that cross-receiver transfer remains unstable.

\begin{figure}[!t]
    \centering
    \includegraphics[width=\linewidth]{fig5_cross_receiver_stress.pdf}
    \caption{Strict source-only cross-receiver stress test. The proposed model improves over CNN-IQ in RX1$\rightarrow$RX2 but not in RX2$\rightarrow$RX1, indicating direction-dependent receiver mismatch rather than receiver-invariant generalization.}
    \label{fig:cross_receiver_stress}
\end{figure}

Overall, the bidirectional average file-level accuracy is 16.7\% for the hybrid model and 13.9\% for CNN-IQ, both far below the 75.0\% cross-day accuracy of the proposed model under the same-receiver setting. Therefore, these results should not be interpreted as evidence of receiver-invariant LoRa RFFI. Instead, they reveal that receiver front-end calibration mismatch remains a major open challenge. The proposed OOB-guided cross-attentive model improves cross-day robustness and some deployment shifts, but strict source-only cross-receiver generalization requires additional receiver-adaptive calibration or multi-receiver training.
